\section{Estimation details}
\label{app:estimation}

\paragraph{Residual construction.} Let $y_t=\operatorname{asinh}(P_t/s_P)$
for hourly PTF from January 2019 to the valuation hour. The model-consistent
residual is $y_t-\bar y_{m(t)}-g_{m(t)}(h(t))$, where $\bar y_m$ is the mean of
$y$ over calendar month $m$ and $g_m$ the mean deviation from it at hour-of-week
$h$ within that month. The construction removes the monthly level, which the
forward curve supplies in the pricing model, and the hourly shape, which the
curve also carries, so that what remains corresponds to $X$ in \eqref{eq:ou}.

\paragraph{Likelihood.} The two-regime autoregression with regime-dependent
mean and volatility and logistic transitions \eqref{eq:tvtp} is estimated by
maximising the likelihood computed with the filter of \citet{Hamilton1989}.
Each fit uses twenty random starting points and a bounded quasi-Newton
optimiser; regime labels are ordered by volatility after estimation, and
standard errors are taken from the inverse Hessian, with outer-product
estimates as a check. The mean-reversion rate is $\kappa=-\ln\phi$ per hour.
Likelihood-ratio statistics for time-varying against constant transitions are
referred to a $\chi^2_2$ distribution.

\paragraph{Boundary estimates.} On the transformed price level without the
monthly level removed, the autoregressive coefficient converges to the unit-root
boundary. Standard asymptotics do not apply there \citep{Andrews2001}, and the
fit is reported through the profile likelihood over a grid of fixed $\phi$,
whose maximum lies at $0.99999$ with a positive definite Hessian for the
remaining parameters. These boundary fits are used only to explain the origin
of the inherited persistence and do not enter pricing.

\paragraph{Real prices.} Deflating PTF by the consumer price index of
TÜİK (series TP.GENENDEKS.T1, base 2003\,=\,100) to December 2025 prices gives
a scale of $3{,}092.05$\,\TRY{} under the documented definition and regime
volatilities 22 to 33\% lower than in nominal terms, so inflation accounts for
part, but not all, of the gap between inherited and re-estimated volatilities.
